%%%%%%%%%%%%%%%%%%%%%%%%%%%%%%%%%%%%%%%%%%%%%%%%%%%%%%%%%%
%   Autoren:
%   Prof. Dr. Bernhard Drabant
%   Prof. Dr. Dennis Pfisterer
%   Prof. Dr. Julian Reichwald
%%%%%%%%%%%%%%%%%%%%%%%%%%%%%%%%%%%%%%%%%%%%%%%%%%%%%%%%%%

%%%%%%%%%%%%%%%%%%%%%%%%%%%%%%%%%%%%%%%%%%%%%%%%%%%%%%%%%%
%	ANLEITUNG: 
%   1. Ersetzen Sie firmenlogo.jpg im Verzeichnis img
%   2. Passen Sie alle Stellen im Dokument an, die mit 
%      @stud 
%      markiert sind
%%%%%%%%%%%%%%%%%%%%%%%%%%%%%%%%%%%%%%%%%%%%%%%%%%%%%%%%%%

%%%%%%%%%%%%%%%%%%%%%%%%%%%%%%%%%%%%%%%%%%%%%%%%%%%%%%%%%%
%	ACHTUNG: 
%   Für das Erstellen des Literaturverzeichnisses wird das 
%   modernere Paket biblatex in Kombination mit biber 
%   verwendet - nicht mehr das ältere Paket BibTex!
%
%   Bitte stellen Sie Ihre TeX-Umgebung entsprechend ein (z.B. TeXStudio): 
%   Einstellungen --> Erzeugen --> Standard Bibliographieprogramm: biber
%%%%%%%%%%%%%%%%%%%%%%%%%%%%%%%%%%%%%%%%%%%%%%%%%%%%%%%%%%

\documentclass[fontsize=12pt,BCOR=5mm,DIV=12,parskip=half,listof=totoc,
               paper=a4,toc=bibliography,pointlessnumbers]{scrreprt}
               
               %toc=listof,listof=entryprefix,
               
\makeindex

%% Elementare Pakete, Konfigurationen und Definitionen werden geladen (gegebenenfalls anpassen)
% !TEX root =  master.tex

%%%%%%%%%%%%%%%%%%%%%%%%%%%%%%%%%%%%%%%%%%%%%%%%%%%%%%%%%%%%%%%%%%
%	ANLEITUNG: 
% Passen Sie gegebenenfalls alle Stellen im Dokument an, die mit 
% @stud 
% markiert sind.
%%%%%%%%%%%%%%%%%%%%%%%%%%%%%%%%%%%%%%%%%%%%%%%%%%%%%%%%%%%%%%%%%%

%%
%% @stud
%%
%% LANGUAGE SETTINGS
\usepackage[ngerman]{babel} 	        % german language
\usepackage[german=quotes]{csquotes} 	% correct quoting using \enquote{}
%\usepackage[english]{babel}          % english language
%\usepackage{csquotes} 	              % correct quoting using \enquote{}
\usepackage{eurosym}
\usepackage{makeidx}                  % allows index generation
\usepackage{listings}	                %Format Listings properly
\usepackage{lipsum}                   % Blindtext
\usepackage{graphicx}       
\usepackage{tabularx}                          % use various graphics formats
\usepackage[german]{varioref}
\usepackage{caption}	                % better Captions
\usepackage{booktabs}                 % nicer Tabs
\usepackage[hidelinks=true]{hyperref} % keine roten Markierungen bei Links
\usepackage{fnpct}                    % Correct superscripts 
\usepackage{calc}                     % Used for extra space below footsepline, in particular
\usepackage{array}
\usepackage{acronym}
\usepackage{algorithm}
\usepackage{algpseudocode}
\usepackage{setspace}
\usepackage{tocloft}
\usepackage{makecell}
\usepackage{hhline}

\usepackage{longtable}
\usepackage{svg}
\usepackage{amssymb}

%Benedikt
\usepackage{xltabular}

% Davids Abschnitt
\usepackage{geometry}
\usepackage{pdflscape}
\usepackage{wrapfig}

%% Schriftarten- und Zeichenpakete
\usepackage[T1]{fontenc}
\usepackage[utf8]{inputenc}

%%
%% @stud
%%
%%	FONT SELECTION: Schriftarten und Schriftfamilie
%%%%%%%%%%%%%
%% SCHRIFTART
%%%%%%%%%%%%%
% 0) without decomment: normal font families 
% ...
% 1) Latin Modern 
\usepackage{lmodern}        
% 2) Times 
%\usepackage{mathptmx}         
% 3) Helvetica
%\usepackage[scaled=.92]{helvet} 
%%%%%%%%%%%%%%%%%%
%%	SCHRIFTFAMILIE
%%%%%%%%%%%%%%%%%%
% ohne Serifen
\renewcommand*{\familydefault}{\sfdefault}
\addtokomafont{disposition}{\sffamily}
%
% mit Serifen
%\renewcommand*{\familydefault}{\rmdefault}
%\addtokomafont{disposition}{\rmfamily}
%
% Typewriter
%\renewcommand*{\familydefault}{\ttdefault}
%\addtokomafont{disposition}{\ttfamily}

%%
%% @stud
%%
%% Uncomment the following lines to support hard URL breaks in bibliography 
%\apptocmd{\UrlBreaks}{\do\f\do\m}{}{}
%\setcounter{biburllcpenalty}{9000}% Kleinbuchstaben
%\setcounter{biburlucpenalty}{9000}% Großbuchstaben

%%
%% @stud
%%
%% FOOTNOTES: Count footnotes over chapters
%% \counterwithout{footnote}{chapter}

%	ACRONYMS
\makeatletter
\@ifpackagelater{acronym}{2015/03/20}
{\renewcommand*{\aclabelfont}[1]{\textbf{{\acsfont{#1}}}}}{}
\makeatother

%	LISTINGS
% @stud: ggf. Namen/Text anpassen (englisch)
\renewcommand{\lstlistingname}{Quelltext} 
\renewcommand{\lstlistlistingname}{Quelltextverzeichnis}
\lstset{numbers=left,
	numberstyle=\tiny,
	captionpos=b,
	basicstyle=\ttfamily\small}

%	ALGORITHMS
% @stud: ggf. Namen/Text anpassen (englisch)
\renewcommand{\listalgorithmname}{Algorithmenverzeichnis}
\floatname{algorithm}{Algorithmus}

%	PAGE HEADER / FOOTER
%	Warning: There are some redefinitions throughout the master.tex-file!  DON'T CHANGE THESE REDEFINITIONS!
\RequirePackage[automark]{scrlayer-scrpage}
%alternatively with separation lines: \RequirePackage[automark,headsepline,footsepline]{scrlayer-scrpage}

\renewcommand{\chaptermarkformat}{}
\RedeclareSectionCommand[beforeskip=0pt]{chapter}
\clearpairofpagestyles

%\ifoot[\rule{0pt}{\ht\strutbox+\dp\strutbox}DHBW Mannheim]{\rule{0pt}{\ht\strutbox+\dp\strutbox}DHBW Mannheim}
\ofoot[\rule{0pt}{\ht\strutbox+\dp\strutbox}\pagemark]{\rule{0pt}{\ht\strutbox+\dp\strutbox}\pagemark}
\ohead{\headmark}

\newcommand{\TitelDerArbeit}[1]{\def\DerTitelDerArbeit{#1}\hypersetup{pdftitle={#1}}}
\newcommand{\AutorDerArbeit}[1]{\def\DerAutorDerArbeit{#1}\hypersetup{pdfauthor={#1}}}
\newcommand{\Firma}[1]{\def\DerNameDerFirma{#1}}
\newcommand{\Kurs}[1]{\def\DieKursbezeichnung{#1}}
\newcommand{\Abteilung}[1]{\def\DerNameDerAbteilung{#1}}
\newcommand{\Studiengangsleiter}[1]{\def\DerStudiengangsleiter{#1}}
\newcommand{\WissBetreuer}[1]{\def\DerWissBetreuer{#1}}
\newcommand{\FirmenBetreuer}[1]{\def\DerFirmenBetreuer{#1}}
\newcommand{\Bearbeitungszeitraum}[1]{\def\DerBearbeitungszeitraum{#1}}
\newcommand{\Abgabedatum}[1]{\def\DasAbgabedatum{#1}}
\newcommand{\Matrikelnummer}[1]{\def\DieMatrikelnummer{#1}}
\newcommand{\Studienrichtung}[1]{\def\DieStudienrichtung{#1}}
\newcommand{\ArtDerArbeit}[1]{\def\DieArtDerArbeit{#1}}
\newcommand{\Literaturverzeichnis}{Literaturverzeichnis}

\newcommand{\DavidsGeileTabellen}[1]{
	\pagebreak
	
	\newgeometry{left=1cm,right=1cm,top=1cm,bottom=1cm}
	\begin{landscape}

		
		\thispagestyle{empty}
		
		\begin{table}
			\begin{tabularx}{27.7cm}{c c p{3cm} X c c c p{2.5cm} c X}
				\textbf{ID} &\textbf{Risikotyp}&\textbf{Titel} &\textbf{Beschreibung} &\textbf{WSK} &\textbf{Stärke} &\textbf{Risiko} &\textbf{Problem Zeitpunkt} &\textbf{MKat} &\textbf{Maßnahme}\\\toprule
				#1
			\end{tabularx}
		\end{table}


	\end{landscape}
	\restoregeometry
}

\newcommand{\settingBibFootnoteCite}{
	\setlength{\bibparsep}{\parskip}		  % Add some space between biblatex entries in the bibliography
	\addbibresource{bibliography.bib}	    % Add file bibliography.bib as biblatex resource
	\DefineBibliographyStrings{ngerman}{andothers = {{et\,al\adddot}},}
}

\newcommand{\setTitlepage}{
	% !TEX root =  master.tex
% @stud: ggf. Namen/Text anpassen (englisch)
\begin{titlepage}
\begin{minipage}{\textwidth}
		\vspace{-2cm}
		\noindent \hfill \includegraphics{\imagedir/logo.jpg}
\end{minipage}
\vspace{1em}
%\sffamily
\begin{center}
	{\textsf{\large Duale Hochschule Baden-W\"urttemberg Mannheim}}\\[4em]
	{\textsf{\textbf{\large{\DieArtDerArbeit}}}}\\[6mm]
	{\textsf{\textbf{\Large{}\DerTitelDerArbeit}}} \\[1.5cm]
	{\textsf{\textbf{\large{}Studiengang Wirtschaftsinformatik}}\\[6mm]
	\textsf{\textbf{Studienrichtung \DieStudienrichtung}}}\vspace{15em}
	
	\begin{minipage}{\textwidth}
		\begin{tabbing}
		Wissenschaftliche(r) Betreuer(in): \hspace{0.85cm}\=\kill
		Verfasser: \> \DerAutorDerArbeit \\[1.5mm]
		Matrikelnummer: \> \DieMatrikelnummer \\[1.5mm]
%		Firma: \> \DerNameDerFirma  \\[1.5mm]
%		Abteilung: \> \DerNameDerAbteilung \\[1.5mm]
		Kurs: \> \DieKursbezeichnung \\[1.5mm]
		Studiengangsleiter: \> \DerStudiengangsleiter \\[1.5mm]
		Dozent: \> \DerWissBetreuer \\[1.5mm]
%		Firmenbetreuer(in): \> \DerFirmenBetreuer \\[1.5mm]
		Bearbeitungszeitraum: \> \DerBearbeitungszeitraum\\[1.5mm]
%		alternativ:\\[1.5mm]
%		Eingereicht: \> \DasAbgabedatum	
		\end{tabbing}
	\end{minipage}
\end{center}
\end{titlepage}
	\pagenumbering{roman} % Römische Seitennummerierung
	\normalfont	
}

\newcommand{\initializeText}{
	\clearpage
	\ihead{\chaptername~\thechapter} % Neue Header-Definition
	\pagenumbering{arabic}           % Arabische Seitenzahlen
}

\newcommand{\initializeBibliography}{
	\ihead{}
	\printbibliography[title=\Literaturverzeichnis] 
	\cleardoublepage
}

\newcommand{\initializeAppendix}{
	\appendix
	\ihead{}
	\cftaddtitleline{toc}{chapter}{Anhang}{}
}


%%
%% @stud
%%
%% PERSÖNLICHE ANGABEN (BITTE VOLLSTÄNDIG EINGEBEN zwischen den Klammern: {...})
%%
\ArtDerArbeit{Projektbericht} % "Bachelor" oder "Projekt" wählen
\TitelDerArbeit{Muli Agent Dev}
\AutorDerArbeit{Benedikt Prisett, David Schäfer,\\[1.5mm] \>Eric Echtermeyer, Lasse Friedrich}
%\Abteilung{<Ihre Abteilung>}
%\Firma{<Ihre Firma>}
\Kurs{WWI-21-DSA}
\Studienrichtung{Data Science}
\Matrikelnummer{5709658, 7086451, 6373947, 9924680}
\Studiengangsleiter{Prof. Dr.-Ing. habil. Dennis Pfisterer}
\WissBetreuer{Clemens Biehl}
%\FirmenBetreuer{<Ihr(e) Firmenbetreuer(in)>}
\Bearbeitungszeitraum{26.11.2023 -- ??.01.2024}

%%
%% @stud
%%
%% BIBLIOGRAPHY (@stud: Bibliographie-Stil wählen - Position und Indizierung)
%%  Auswahl zwischen: 
%%   NUMERIC Style
%%   IEEE Style
%%   ALPHABETIC Style
%%   HARVARD Style
%%   CHICAGO Style 
%%   (oder eigenen zulässigen Stil wählen) 
%%
%%%%%%%%%%%%%
%% Zitierstil
%%%%%%%%%%%%%
% NUMERIC Style - e. g. [12]
\newcommand{\indextype}{numeric} 
%
% IEEE Style - numeric kind of style 
%\newcommand{\indextype}{ieee} 
%
% ALPHABETIC Style - e. g. [AB12]
%\newcommand{\indextype}{alphabetic} 
%
% HARVARD Style 
%\newcommand{\indextype}{apa} 
%
% CHICAGO Style 
%\newcommand{\indextype}{authoryear}
%
%%%%%%%%%%%%%%%%%%%%%%
%% Position des Zitats
%%%%%%%%%%%%%%%%%%%%%%
\newcommand{\position}{inline} 
%
% (!!) FOOTNOTE POSITION NOT RECOMMENDED IN MINT DOMAIN:
%\newcommand{\position}{footnote}

%% Final: Setzen des Zitierstils und der Zitatposition
\usepackage[backend=biber, autocite=\position, style=\indextype]{biblatex} 	
\settingBibFootnoteCite

%%
%% Definitionen und Commands
%%
\newcommand{\abs}{\par\vskip 0.2cm\goodbreak\noindent}
\newcommand{\nl}{\par\noindent}
\newcommand{\mcl}[1]{\mathcal{#1}}
\newcommand{\nowrite}[1]{}
\newcommand{\NN}{{\mathbb N}}
\newcommand{\imagedir}{img}

\makeindex

\begin{document}

\setTitlepage

%%%%%%%%%%%%%%%%%%%%%%%%%%%%%%%%%%%
% EHRENWÖRTLICHE ERKLÄRUNG
%
% @stud: ewerkl.tex bearbeiten
%
%% !TEX root =  master.tex
\clearpage
\chapter*{Ehrenwörtliche Erklärung}

% Wird die folgende Zeile auskommentiert, erscheint die ehrenwörtliche
% Erklärung im Inhaltsverzeichnis.

% \addcontentsline{toc}{chapter}{Ehrenwörtliche Erklärung}
Ich versichere hiermit, dass ich die vorliegende Seminararbeit mit dem Titel ``\textit{\DerTitelDerArbeit}'' selbstständig verfasst und 
keine anderen als die angegebenen Quellen und Hilfsmittel benutzt habe. Ich versichere zudem, dass die eingereichte elektronische 
Fassung mit der gedruckten Fassung übereinstimmt.

\vspace{3cm}
Ort, Datum \hfill \DerAutorDerArbeit
 
%\cleardoublepage  
%%%%%%%%%%%%%%%%%%%%%%%%%%%%%%%%%%%

%%%%%%%%%%%%%%%%%%%%%%%%%%%%%%%%%%%
% SPERRVERMERK
%
% @stud: nondisclosurenotice.tex bearbeiten
%
%% !TEX root =  master.tex
\chapter*{Sperrvermerk}

\begin{center}
\fbox{
		\begin{minipage}{33em}
			\textbf{Ein Sperrvermerk sollte nur bei berechtigtem Bedarf gesetzt werden!\\[10pt] 
				Beachten Sie, dass mit Sperrvermerk	versehene Arbeiten nicht für weitere wissenschaftliche Zwecke 
				außerhalb des Firmenkontextes oder zur Publikation verwendet werden dürfen.\\[10pt]
				Wir empfehlen, wenn m\"oglich, auf den Sperrvermerk zu verzichten.\\[10pt]
				Besprechen Sie diese Problematik mit Ihrer Firma!}
		\end{minipage}
}
\end{center}

(Mustertext) Der Inhalt dieser Arbeit darf weder als Ganzes noch in Auszügen Personen außerhalb des Prüfungsprozesses 
und des Evaluationsverfahrens zugänglich gemacht werden, sofern keine anders lautende Genehmigung der Ausbildungsstätte vorliegt. 

\cleardoublepage
 
%\cleardoublepage
%%%%%%%%%%%%%%%%%%%%%%%%%%%%%%%%%%%

%%%%%%%%%%%%%%%%%%%%%%%%%%%%%%%%%%%
%	KURZFASSUNG
%
% @stud: acknowledge.tex bearbeiten
%
%% !TEX root =  master.tex
\chapter*{Danksagung}

Hier können Sie eine Danksagung schreiben. 



%\cleardoublepage 
%%%%%%%%%%%%%%%%%%%%%%%%%%%%%%%%%%%

%%%%%%%%%%%%%%%%%%%%%%%%%%%%%%%%%%%
% VERZEICHNISSE und ABSTRACT
%
% @stud: ggf. nicht benötigte Verzeichnisse auskommentieren/löschen
%
\tableofcontents
\cleardoublepage

% Abbildungsverzeichnis
\phantomsection
\addcontentsline{toc}{chapter}{\listfigurename}
\listoffigures
\cleardoublepage

%	Tabellenverzeichnis
\phantomsection
\addcontentsline{toc}{chapter}{\listtablename}
\listoftables
\cleardoublepage

%	Listingsverzeichnis / Quelltextverzeichnis
%\lstlistoflistings
%\cleardoublepage

% Algorithmenverzeichnis
%\listofalgorithms
%\cleardoublepage

% Abkürzungsverzeichnis
% @stud: acronyms.tex bearbeiten
%% !TEX root =  master.tex
\clearpage
\chapter*{Abkürzungsverzeichnis}	
\addcontentsline{toc}{chapter}{Abkürzungsverzeichnis}

\begin{acronym}[XXXXXXX]
	\acro{ad}[AD]{Archiv für Diplomatik, Schriftgeschichte, Siegel- und Wappenkunde}
	\acro{BMBF}{Bundesministerium für Bildung und Forschung}	
	\acro{DHBW}{Duale Hochschule Baden-Württemberg}
	\acro{ecu}[ECU]{European Currency Unit}
	\acro{eu}[EU]{Europäische Union}
	\acro{RDBMS}[RDBMS]{Relational Database Management System}
\end{acronym} 
%\cleardoublepage

%	Kurzfassung / Abstract
% @stud: abstract.tex bearbeiten
% !TEX root =  master.tex
%%
%% @stud: sprachspezifische Anpassung
%%
\chapter*{Kurzfassung (Abstract)}
\addcontentsline{toc}{chapter}{Kurzfassung (Abstract)} 
\cleardoublepage

%%%%%%%%%%%%%%%%%%%%%%%%%%%%%%%%%%%%%%%%%%%%%%%%%%%%%%%%%%%%%%%%%%%%%%%%%%%%%%%%%%%%%%%%%%
% KAPITEL UND ANHÄNGE
%
% @stud:
%   - nicht benötigte: auskommentieren/löschen
%   - neue: bei Bedarf hinzufügen mittels input-Komman do an entsprechender Stelle einfügen
%%%%%%%%%%%%%%%%%%%%%%%%%%%%%%%%%%%%%%%%%%%%%%%%%%%%%%%%%%%%%%%%%%%%%%%%%%%%%%%%%%%%%%%%%%

\initializeText
\onehalfspacing

%%%%%%%%%%%%%%%%%%%%%%%%%%%%%%%%%%%
% KAPITEL
%
% @stud: einzelne Kapitel bearbeiten und eigene Kapitel hier einfügen
%
% Einleitung
%% !TEX root =  master.tex
\chapter{Einleitung}

\nocite{*}

Dieses Kapitel enthält die Einleitung mit ihren verschiedenen Abschnitten/Sections und Unterabschnitten.

\section{Beispiel Abschnitt: \LaTeX-Installation}

Zur Verwendung von \LaTeX-Installation einer Distribution z.~B.~TeXLive, MikTex etc.~sowie eines Editors z.~B.~TeXStudio, TeXnicCenter etc.~notwendig.

Installieren Sie zun\"achst die Distribution und anschließend den Editor. Beim ersten Start des Editors \"offnet sich ein 
Konfigurationsassistent, der zun\"achst nach dem Pfad der installierten Distribution fragt. 

Nach der Installation können k\"onnen Einstellungen z.~B.~f\"ur einen PostScript-Viewer gemacht werden. 
Dieser Schritt kann ohne Weiteres \"ubersprungen werden. Entscheidend sind die Einstellungen f\"ur den pdf-Viewer. 

Jetzt kann \LaTeX~verwendet werden. Um die Ausgabe eines Dokumentes zu erzeugen, muss das Dokument kompiliert werden (Ausgabe >
Aktives Dokument > Erstellen und betrachten).

\subsection{Beispiel Unterabschnitt: Aufbau eines \LaTeX-Dokuments}

Ein \LaTeX-Dokument besteht in der Regel aus folgenden Komponenten:
\begin{itemize}
	\item Pr\"aambel
	\item Titelseite
	\item Textteil
\end{itemize}

\subsection{Beispiel Unterabschnitt auf zweiter Ebene: Pr\"aambel}
In der Pr\"aambel werden global die Einstellungen f\"ur das gesamte Dokument definiert. Hierbei k\"onnen z.~B.~die Seitenr\"ander, 
der Zeilenabstand oder auch die Sprache f\"ur die Silbentrennung festgelegt werden. In der ersten Zeile eines jeden Dokumentes wird dabei
immer die zu verwendende Klasse festgelegt. Standardm\"aßig kann hier die Artikel-Klasse gew\"ahlt werden:

\texttt{\textbackslash documentclass[12pt,titlepage]\{article\}}

In den eckigen Klammern wird dabei u.a. die Standardschriftgr\"o\ss e f\"ur das gesamte Dokument festgelegt. 

Au\ss erdem werden in der Pr\"aambel die f\"ur das Dokument ben\"otigten Pakete festgelegt. Gebr\"auchlich sind vor allem folgende Pakete:
{\texttt{
\begin{itemize}
	\item \textbackslash usepackage[ngerman]\{babel\}
	\item \textbackslash usepackage[latin1]\{inputenc\}
	\item \textbackslash usepackage\{color\}
	\item \textbackslash usepackage[a4paper]\{geometry\}
	\item \textbackslash usepackage\{amssymb\}
	\item \textbackslash usepackage\{amsthm\}
	\item \textbackslash usepackage\{graphicx\}
\end{itemize}
}

Im vorliegenden Fall werden die Pakete in der Konfigurationsdatei \texttt{config.tex} festgelegt, deren Inhalt durch 
\texttt{\textbackslash input\{config\}} in das Hauptdokument \texttt{master.tex} inkludiert wird.

\subsubsection{Beispiel Unterabschnitt auf zweiter Ebene: Titelseite}

Nachdem die Dokumenten-Klasse und die zu verwendenden Pakete festgelegt worden sind,
folgt die Titelseite. Da die Titelseite bereits Teil des eigentlichen Dokuments ist, muss ihr
unbedingt der Befehl \texttt{\textbackslash begin\{document\}} vorausgehen. Am Ende des Dokuments sollte der Befehl
\texttt{\textbackslash end\{document\}} gesetzt werden. Alles was nach diesem Befehl steht, wird vom Compiler nicht mehr beachtet.

\section{Noch ein Beispiel-Abschnitt}

Der Textteil beinhaltet nun den eigentlichen Text des Dokuments.


% mehrere Grundlagen- und Forschungs-Kapitel
% !TEX root = master.tex
\chapter{Einleitung}



\chapter{Aufbau}

% !TEX root = master.tex
\chapter{Auswertung}

% !TEX root = master.tex
\chapter{Chapter 4}

% !TEX root = master.tex
\chapter{Chapter 5}

% !TEX root = master.tex
\chapter{Chapter 6}

% Fazit und Ausblick
%% !TEX root =  master.tex
\chapter{Zusammenfassung}

\nocite{*}

Dieses Kapitel enthält die Zusammenfassung der Arbeit mit Fazit und Ausblick.

\section{Fazit}

...

\section{Ausblick}

...


%%%%%%%%%%%%%%%%%%%%%%%%%%%%%%%%%%%

%%%%%%%%%%%%%%%%%%%%%%%%%%%%%%%%%%%
% ANHÄNGE
%
% @stud: einzelne Anhänge bearbeiten und eigene Anhänge hier einfügen 
%        die nachfolgenden Zeilen deaktivieren, wenn keine Anhänge verwendet werden
%
%\initializeAppendix
%% !TEX root =  master.tex
\chapter{Beispiel-Anhang: Testanhang}

%% !TEX root =  master.tex
\chapter{Beispiel-Anhang: Noch ein Testanhang}
nochmal: lipsum ...

%%%%%%%%%%%%%%%%%%%%%%%%%%%%%%%%%%%

\singlespacing

%%%%%%%%%%%%%%%%%%%%%%%%%%%%%%%%%%%
% LITERATURVERZEICHNIS
% @stud: Literaturverzeichnis in Datei bibliography.bib anpassen. 
%
% Alternative zu Verwendung von \initializeBibliography: Citavi ...
% (dann \initializeBibliography auskommentieren und eigenes LaTex Coding verwenden)
%
\initializeBibliography
%%%%%%%%%%%%%%%%%%%%%%%%%%%%%%%%%%%

%%%%%%%%%%%%%%%%%%%%%%%%%%%%%%%%%%%
% INDEX
% @stud: ggf. Index auskommentieren, wenn nicht benötigt
%
%\addcontentsline{toc}{chapter}{Index}
%\printindex

\end{document}
